%=============================================================================
% APPENDIX S: STANDARD MODEL EXTENSION DICTIONARY
% DFD Unified Review v3.0
%=============================================================================

\section{Standard Model Extension Dictionary}\label{app:sme}

This appendix maps \DFD parameters onto the language of the Standard-Model 
Extension (SME)~\cite{Kostelecky2011}, enabling direct comparison with 
published experimental constraints.

%-----------------------------------------------------------------------------
\subsection{SME Framework Overview}\label{app:sme-overview}
%-----------------------------------------------------------------------------

The SME provides a phenomenological framework for parameterizing possible 
violations of Lorentz invariance and the Einstein Equivalence Principle. 
For gravitational tests with atomic clocks, the relevant observable is:
\begin{equation}\label{eq:sme-observable}
\frac{\delta(f_A/f_B)}{(f_A/f_B)} = (\beta_A - \beta_B)\frac{\Delta U}{c^2},
\end{equation}
where $\beta_A, \beta_B$ encode gravitational redshift anomalies for species $A$ and $B$.

%-----------------------------------------------------------------------------
\subsection{DFD$\leftrightarrow$SME Correspondence}\label{app:sme-mapping}
%-----------------------------------------------------------------------------

In \DFD, the same observable is:
\begin{equation}\label{eq:dfd-clock-obs}
\frac{\delta(f_A/f_B)}{(f_A/f_B)} = (\xi_A - \xi_B)\frac{\Delta\Phi}{c^2},
\end{equation}
where the effective coupling $\xi_A$ includes both matter and photon sector contributions:
\begin{equation}\label{eq:xi-def}
\xi_A \equiv K_A + \delta_{A,\gamma},
\end{equation}
with the full channel-resolved coupling $K_A$ from Eq.~\eqref{eq:K-general} (of which the pure-$\alpha$ leading term is $K_A^{(\alpha)} = \kalpha \cdot S_A^\alpha$) and 
$\delta_{A,\gamma} = 1$ if species $A$ involves a photon-sector reference.
After the constitutive-chain cancellation of Sec.~\ref{sec:cavity}, the photon-sector contribution $\delta_{A,\gamma}$ is absorbed into the tree-level cancellation and the surviving observable is a screened residual.

Identifying $\Delta U \leftrightarrow \Delta\Phi$, the direct correspondence is:
\begin{equation}\label{eq:sme-dfd-map}
\boxed{\beta_A - \beta_B \longleftrightarrow \xi_A - \xi_B = (K_A - K_B) + (\delta_{A,\gamma} - \delta_{B,\gamma})}
\end{equation}

%-----------------------------------------------------------------------------
\subsection{Translation Table}\label{app:sme-table}
%-----------------------------------------------------------------------------

\begin{table}[h]
\centering
\caption{DFD$\leftrightarrow$SME parameter correspondence.}
\label{tab:sme-dfd}
\scriptsize
\setlength{\tabcolsep}{2pt}
\begin{tabular}{@{}lll@{}}
\toprule
\textbf{DFD Quantity} & \textbf{SME Analogue} & \textbf{Meaning} \\
\midrule
$\psi$ & $U/c^2$ & Background grav.\ field \\
$K_A$ & $\beta_A$ (matter) & Species-dep.\ coupling \\
$\delta_{A,\gamma}$ & $\beta_A$ (photon) & Photon-mode coupling \\
$\xi_A = K_A + \delta_{A,\gamma}$ & Total $\beta_A$ & Composite LPI param. \\
$\kalpha = \alpha^2/(2\pi)$ & --- & \DFD-specific scale \\
\bottomrule
\end{tabular}
\end{table}

%-----------------------------------------------------------------------------
\subsection{Experimental Constraints Reinterpreted}\label{app:sme-experimental}
%-----------------------------------------------------------------------------

Published SME bounds can be reinterpreted as \DFD constraints:

\begin{table}[h]
\centering
\caption{SME bounds reinterpreted in \DFD framework.}
\label{tab:sme-bounds}
\resizebox{\columnwidth}{!}{%
\begin{tabular}{@{}lccl@{}}
\toprule
\textbf{Experiment} & \textbf{SME Constraint} & \textbf{DFD Interpretation} & \textbf{Ref.} \\
\midrule
H maser/Cs (14-yr) & $|\beta_{\rm H} - \beta_{\rm Cs}| < 2.5 \times 10^{-7}$ & $|K_{\rm H} - K_{\rm Cs}| < 2.5 \times 10^{-7}$ & \cite{Ashby2018} \\
Yb$^+$ E3/E2 (PTB) & $|\beta_{\rm E3} - \beta_{\rm E2}| < 10^{-8}$ & Same-ion: composition cancels & \cite{Filzinger2023} \\
Hg$^+$/Cs & $|\beta_{\rm Hg} - \beta_{\rm Cs}| < 5.8 \times 10^{-6}$ & $|K_{\rm Hg} - K_{\rm Cs}| < 5.8 \times 10^{-6}$ & \cite{Fortier2007} \\
Al$^+$/Hg$^+$ & $|\beta_{\rm Al} - \beta_{\rm Hg}| < 5.3 \times 10^{-7}$ & $|K_{\rm Al} - K_{\rm Hg}| < 5.3 \times 10^{-7}$ & \cite{Rosenband2008} \\
\bottomrule
\end{tabular}
}% end resizebox
\end{table}

%-----------------------------------------------------------------------------
\subsection{Cavity-Atom Comparisons in SME Language}\label{app:sme-cavity}
%-----------------------------------------------------------------------------

For cavity-atom LPI tests (Sec.~\ref{sec:cavity}), the SME parameterization becomes:
\begin{equation}
\frac{d}{d\Phi}\left(\frac{\nu_{\rm atom}}{\nu_{\rm cavity}}\right) = \frac{\xi_{\rm atom} - \xi_{\rm cavity}}{c^2} = \frac{\xi_{\rm LPI}}{c^2},
\end{equation}
where the old tree-level assignment $\xi_{\rm cavity}=1$ is no longer used. After the constitutive-chain cancellation of Sec.~\ref{sec:cavity}, both cavity and atomic sectors share the universal geometric redshift at tree level, so only a screened residual mismatch survives:
\begin{equation}
\xi_{\rm LPI}\ \longrightarrow\ \xi_{\rm LPI}^{\rm res}.
\end{equation}

\paragraph{Significance.}
In SME-style language, \DFD no longer predicts a dramatic order-unity cavity coefficient. Instead it predicts that any measurable cavity--atom anomaly must arise from a channel-resolved residual, consistent with the four-term structure of Eq.~\eqref{eq:K-general} and the screening logic summarized in Secs.~\ref{sec:clocks} and \ref{sec:cavity}.
